\makeatletter
% Add the ICLR 2026 template directory to LaTeX's search path.
% This lets \usepackage{iclr2026_conference} work without copying files.
\def\input@path{{../doc/iclr-template/master/iclr2026/}}
\makeatother

\documentclass{article}

% ICLR template (located in ../doc/templates/iclr2026/)
\usepackage{iclr2026_conference}
\usepackage{times}

% Math & symbols
\usepackage{amsmath,amsfonts,amssymb,amsthm}

% Figures & tables
\usepackage{graphicx}
\usepackage{booktabs}
\usepackage{subcaption}

% Formatting
\usepackage{microtype}
\usepackage{hyperref}
\usepackage{xcolor}
\usepackage{enumitem}
\newcommand{\draft}[1]{}

% Reduce underfull vbox warnings by allowing variable page height.
\raggedbottom

% Theorem-like environments (used in \input files).
\newtheorem{theorem}{Theorem}
\newtheorem{lemma}{Lemma}
\newtheorem{proposition}{Proposition}
\newtheorem{corollary}{Corollary}
\newtheorem{assumption}{Assumption}

\title{The Curse and Blessing of Bias in Large Language Model Training}

\author{
  Anonymous Authors\\
  Anonymous Institutions\\
  \texttt{anonymous@iclr.cc}
}

% Uncomment for camera-ready
% \iclrfinalcopy

\begin{document}

\maketitle

\begin{abstract}

Human language corpora, the fuel of intelligence for large language models (LLMs), induce a strongly anisotropic geometry in learned representations.
Across layers and training stages, activations and parameters exhibit a pronounced spike--tail structure:
a small set of directions concentrates disproportionate energy,
while the remaining directions form a broad, information-rich long tail.
Such anisotropy acts as a \emph{curse} for large-scale training.
High-energy components inflate elementwise dynamic range,
distort adaptive second-moment estimates,
and compress low-energy semantic variation in low-bit regimes.

However, a closer examination reveals that anisotropy is also a \emph{blessing}.
The dominant spike is not arbitrary—it is highly coherent and largely driven by a shared mean bias that originates in the embedding layer and propagates through residual connections.
Because this bias is structured and low-rank,
it admits an exceptionally simple statistical description.
In particular, a single averaging operation can isolate and remove the dominant mean component without explicit spectral decomposition.
This operation dramatically reduces elementwise outliers and second-moment inflation,
while preserving the semantic covariance carried by the long tail.

In this work, we formalize this dual role of anisotropy and show that
a simple bias-centric spectral conditioning strategy—
mean removal, optionally combined with lightweight spike whitening—
can \emph{replace or complement} SVD-based spike suppression.
We provide theoretical analysis explaining why mean bias dominates extreme behavior,
and we empirically validate the approach in large language model training,
including downstream evaluations in low-bit quantization and adaptive optimization.

\end{abstract}


\section{Introduction}

Large language models are trained on massive human language corpora,
which induce a characteristic geometry in learned representations.
Across layers and training stages, activations, gradients, and model parameters
exhibit pronounced anisotropy:
a small number of directions concentrate disproportionate energy
(strong spectral spikes),
while the remaining directions form a broad long tail
\cite{ethayarajh2019contextual,mu2018all,timkey2021rogue,kovaleva2021bertbusters,machina2024anisotropy}.

Although such structure reflects statistical regularities of language,
shared grammatical patterns versus diverse semantic content, it creates computational challenges.
Anisotropy inflates elementwise dynamic range,
distorts adaptive second-moment estimates \cite{kingma2014adam,loshchilov2017adamw,spectra2026},
and compresses long-tail variation in low-bit regimes
\cite{xiao2023smoothquant,frantar2022gptq,lin2024awq,li2024svdquant,ashkboos2024quarot}.
Consequently, training becomes numerically unstable or inefficient
unless spectral structure is explicitly controlled.

In this work, we investigate the origin of anisotropy across activations, gradient outputs, and model parameters. We uncover a structural asymmetry: \emph{anisotropy in activations is dominated by a coherent mean bias, whereas gradient outputs exhibit no comparable mean component}. Tracing this structure backward, we identify its source in the embedding layer. Because token frequencies in natural language are highly skewed, embedding weights develop a frequency-aligned bias \cite{mikolov2013efficient,mikolov2013distributed,arora2017simple}. Through residual connections, this bias propagates across layers, inducing persistent anisotropy throughout the network.

This reframes anisotropy at its root.
Rather than being a diffuse high-dimensional artifact,
it is largely anchored in a structured, frequency-dominated mean component.
Crucially, this structure is not merely harmful—it is also exploitable.
Because the dominant instability arises from a rank-one bias,
it can be isolated and removed with a single mean-subtraction operation,
consistent with prior embedding post-processing approaches
\cite{arora2017simple,mu2018all,huang2021whiteningbert}.
The remaining low-rank spike can, if necessary,
be whitened once at the embedding activation.
Both operations are hardware-friendly and substantially cheaper
than repeated layer-wise SVD corrections.
In this sense, the same structure that creates instability
also enables efficient mitigation—the \emph{blessing} of bias.

We validate this bias-centric anisotrophy in two downstream settings.
In low-bit quantization, mean removal significantly reduces dynamic-range inflation
\cite{xiao2023smoothquant,wei2023outlier,ashkboos2024quarot}.
In adaptive optimization, it contracts extreme second-moment tails,
improving effective learning-rate allocation
\cite{kingma2014adam,spectra2026}.
These results suggest that addressing bias at its source
can replace or complement spectral interventions
with minimal computational overhead.


\paragraph{Contributions.}
\begin{enumerate}[leftmargin=*,itemsep=2pt]
    \item We identify activation mean bias as the primary structural source of anisotropy in large language model training.
    \item We show that mean-dominated components account for most extreme activation magnitudes and second-moment inflation.
    \item We trace this bias to skewed token frequencies in the embedding layer and its propagation through residual connections.
    \item We propose source-level mean removal and spike whitening as efficient spectral conditioning mechanisms.
    \item We demonstrate improvements in low-bit quantization and adaptive optimization.
\end{enumerate}


















\draft{
Large language models are trained on massive human language corpora.
Across architectures and scales, these corpora induce a characteristic geometry in
learned representations.
Empirically, activations and gradients in modern language models are highly
anisotropic: a small number of directions exhibit unusually large energy,
forming strong spectral spikes, while the remaining directions constitute a
broad, information-rich long tail.
This spike--tail structure appears consistently across layers, models, and
training regimes.

This anisotropy is not an artifact of a particular model or optimizer.
Rather, it is a fundamental property of human language data.
Linguistic structure---grammar, syntax, and shared functional patterns---is
highly coherent across tokens and sentences.
In contrast, semantic content is diverse, context-dependent, and weakly
correlated across tokens.
When embedded into a high-dimensional representation space, this statistical
asymmetry manifests as a low-rank linguistic subspace that induces spectral
spikes, alongside a high-dimensional semantic subspace that forms a long tail.

This structure is a \emph{curse}.
Strong anisotropy leads to poorly conditioned optimization landscapes, suppresses
learning in semantically meaningful tail directions, and interacts unfavorably
with adaptive optimizers such as Adam and AdamW.
In low-bit or mixed-precision settings, the dominant spike controls numerical
scale and clipping behavior, causing severe precision loss in the long tail.
As a result, aggressive quantization and large-scale optimization become
unstable or inefficient unless anisotropy is explicitly addressed.

The standard approach to controlling anisotropy is spectral.
Methods based on singular value decomposition (SVD) or related techniques
explicitly identify and suppress the dominant spike subspace, and have proven
effective in both optimization and quantization settings.
However, these approaches are computationally expensive and hardware unfriendly,
requiring QR decompositions or iterative orthogonalization that do not map well
to modern accelerators.
This limits their applicability in large-scale training loops.

In this work, we identify a complementary and previously underexplored aspect of
anisotropy: \emph{anisotropy is not only a curse, but also a blessing}.
The same statistical regularities that give rise to spectral spikes imply that
linguistic structure is coherent and shared across tokens, while semantic
components are diverse and weakly correlated under batch sampling. In other words, linguistic structure is the shared frozen component, while semantics appears primarily as higher-order structured variability captured by covariates.
As a consequence, the dominant spike admits a simple statistical
characterization: a substantial portion of its energy concentrates in shared,
low-dimensional components that can be accessed without explicit spectral
decomposition.

Leveraging this observation, we show that a much simpler operation---\emph{averaging}---can
serve as a principled and efficient spectral approximation.
By subtracting the batch mean, we remove a shared rank-one component that lies
primarily within the low-rank linguistic spike subspace, while preserving the
semantic covariance that encodes learnable information.
Unlike SVD-based spike removal, averaging is conservative: it suppresses the most
harmful shared structure without destroying semantic variation, yet it can
\emph{fill in the role of SVD-based spike suppression} in practice.

We formalize this observation through a simple model that separates linguistic
and semantic components, provide a theoretical analysis showing that averaging
removes shared linguistic structure while preserving semantic covariance up to
vanishing error, and empirically validate the theory using activations from large
language models.

We demonstrate the practical impact of this perspective in two complementary
settings.
First, in low-bit quantization, averaging reduces spike-dominated scale effects
and preserves long-tail resolution, explaining and complementing recent spectral
methods for FP4 training.
Second, in optimization, averaging attenuates spike-dominated second-moment
statistics that suppress tail learning under adaptive optimizers, providing a
lightweight approximation to spike-aware optimization techniques.
In both cases, averaging achieves the desired spectral control while avoiding
costly decompositions.

Taken together, this work highlights a dual role of anisotropy in human language
corpora.
While anisotropy introduces fundamental challenges for optimization and
quantization, it also enables efficient remedies that exploit the coherence of
linguistic structure.
By aligning algorithmic design with data geometry, we obtain a simple,
theoretically grounded, and hardware-friendly approach to mitigating the
pathologies induced by anisotropic language representations.

\paragraph{Contributions.}
This work makes the following contributions:
\begin{itemize}
    \item We characterize anisotropy as a fundamental geometric property of human
    language corpora, arising from the coexistence of shared linguistic structure and
    diverse semantic content, and show how this structure manifests as a low-rank
    spectral spike with a high-dimensional semantic tail in large language model
    representations.
    \item We formalize the dual role of anisotropy as both a \emph{curse} and a
    \emph{blessing}: while the spike--tail structure can degrade optimization and
    quantization by suppressing long-tail directions, the same statistical
    regularities imply that the dominant linguistic component is coherent and shared
    across tokens, enabling efficient isolation.
    \item We develop a theoretical framework explaining why subtracting the batch mean
    removes shared linguistic structure while preserving semantic covariance up to
    vanishing error, clarifying where information relevant to learning resides.
    \item We introduce averaging-based spectral approximation as a principled,
    hardware-friendly alternative that can fill in the role of SVD-based spike
    suppression, and empirically validate its effectiveness on large language model
    activations by explaining and complementing recent advances in low-bit
    quantization and spike-aware optimization.
\end{itemize}
}

% \section{Empirical Analysis of Anisotropy in Language Model Representations}

The purpose of this section is to empirically characterize the geometric structure
induced by human language corpora in large language model representations.
Rather than proposing new algorithms, we focus on establishing a set of
data-driven facts about anisotropy that hold across models, layers, and training
settings.
In particular, we show that anisotropy consistently manifests as a separation
between a coherent, low-rank spectral spike and an incoherent, high-dimensional
long tail.
This empirical characterization provides the foundation for the averaging-based
spectral approximation developed in subsequent sections.

\subsection{Experimental Setup}

Our analysis is based on activations extracted from pretrained large language
models.
Unless otherwise stated, we analyze intermediate activation tensors from a
decoder-only transformer, using a batch of sequences drawn from a standard
language modeling corpus.
For a given layer, activations have shape $[B, L, H]$, where $B$ is the batch
size, $L$ is the sequence length, and $H$ is the hidden dimension.
We flatten the batch and sequence dimensions to obtain a token matrix
$X \in \mathbb{R}^{N \times H}$ with $N = B \cdot L$ tokens.

All analyses in this section are performed on raw activations without
preprocessing or normalization unless explicitly stated.
We emphasize that the observations reported below are qualitatively stable across
layers, batches, and model scales, and do not depend on architectural details or
training objectives.

\subsection{Spectral Signature: Spike and Long Tail}
\label{subsec:spike_long_tail}

We begin by examining the singular value spectrum of the token matrix
$b \in \mathbb{R}^{(BL)\times H}$ obtained by flattening the batch and sequence
dimensions of the activation tensor.
Let
\[
b = U \Sigma V^\top,
\qquad
\Sigma = \mathrm{diag}(\sigma_1,\sigma_2,\dots),
\]
denote its singular value decomposition.
Throughout this section, we define \emph{spectral energy} as the squared singular
values $\sigma_k^2$, so that the total energy equals the Frobenius norm
$\sum_k \sigma_k^2 = \|b\|_F^2$.

Figure~\ref{fig:spectrum_spike_tail_combined} summarizes the spectral signature
of anisotropy using the singular value spectrum and the cumulative spectral
energy distribution.
Across all layers and batches we analyzed, the spectrum exhibits a striking and
stable spike--tail structure: a small number of singular values dominate the
spectrum, followed by a broad and smoothly decaying long tail.

\begin{figure}[t]
  \centering
  \begin{subfigure}[t]{0.48\linewidth}
    \centering
    \includegraphics[width=\linewidth]{figures/spectrum_spike_tail.png}
    \caption{Singular value spectrum (log scale).}
    \label{fig:spectrum_spike_tail}
  \end{subfigure}
  \hfill
  \begin{subfigure}[t]{0.48\linewidth}
    \centering
    \includegraphics[width=\linewidth]{figures/spectrum_cum_energy.png}
    \caption{Cumulative spectral energy
    $\sum_{i\le k}\sigma_i^2 / \sum_i\sigma_i^2$.}
    \label{fig:spectrum_cum_energy}
  \end{subfigure}

  \caption{Spectral signature of anisotropy in language model activations.
  (\textbf{Left}) The singular value spectrum exhibits a sharp low-rank spike
  followed by a broad and smoothly decaying long tail.
  (\textbf{Right}) The cumulative spectral energy shows that a very small fraction
  of components captures a large share of the total energy, while hundreds to
  thousands of directions are required to explain the remaining variance,
  indicating a broad and information-rich tail.}
  \label{fig:spectrum_spike_tail_combined}
\end{figure}

Quantitatively, in the example shown, the top $1\%$ of singular components capture
more than $60\%$ of the total spectral energy.
Moreover, fewer than $20$ components suffice to explain $50\%$ of the energy,
while reaching $90\%$ and $95\%$ requires several hundred components
(e.g., $k\approx 390$ and $k\approx 799$, respectively).
These values are computed directly from the cumulative energy distribution shown
in Figure~\ref{fig:spectrum_cum_energy}, and are explicitly annotated in the plot.
At the same time, the singular value spectrum continues to decay smoothly over
thousands of directions, indicating that the long tail is broad and structured,
rather than negligible noise.

This spike--tail separation is robust across layers, model scales, and training
regimes, and has been independently observed in prior studies on low-bit
quantization and optimizer dynamics.
From a data-geometry perspective, these results establish anisotropy as a stable
and intrinsic property of language model representations.
They also motivate a key design principle for mitigating anisotropy: any
effective intervention must selectively suppress the dominant spike while
preserving the high-dimensional structure of the long tail, which encodes
semantically meaningful variation.

While this section establishes the existence and scale of the spectral spike, it
does not address whether this spike corresponds to shared structure or random
variance.
We address this question next.


\subsection{The Spike: Low-Rank and Coherent}
\label{subsec:spike_coherent}

We now examine the structure of the dominant spectral spike and show that it is
both low-rank and highly coherent across tokens.
This distinguishes the spike from random high-variance noise and provides the
key empirical foundation for averaging-based spike suppression.

\paragraph{Low-rank structure of the spike.}
Although the leading singular value is often emphasized, the spike is not
confined to a single direction.
Instead, its energy is distributed across a small number of leading singular
components, forming a low-dimensional subspace.
In the example analyzed here, the top singular component captures approximately
$32.5\%$ of the batch mean energy, while the top $5$ components already capture
more than $96\%$.
Projecting onto the top $10$ components captures over $99\%$ of the batch mean
energy, and projecting onto the top $1\%$ of singular directions ($k=41$)
recovers more than $99.9\%$.
This demonstrates that the spike is low-rank but not rank-one, and that the
shared structure is distributed across a small subspace rather than a single
singular direction.

\paragraph{Coherence with shared structure.}
Large singular values alone do not imply shared structure: a high-variance
direction could arise from uncorrelated fluctuations.
To directly test whether the spike corresponds to a coherent, token-shared
component, we compare the broadcasted column-mean matrix
$\mathrm{ColMean}_{\mu}(b)=\mathbf{1}\mu^\top$ with individual rank-$1$
reconstructions $\sigma_k u_k v_k^\top$ of the activation matrix $b$.
For each singular component $k$, we measure the absolute cosine similarity in
Frobenius inner product, which captures how strongly the shared mean aligns with
that component.

Figure~\ref{fig:colmean_alignment} shows the resulting alignment curves for three
different estimators of the column mean.
The horizontal axis is plotted on a logarithmic scale to emphasize the sharp
localization of alignment in the leading components.
Across all estimators, alignment is highly concentrated in the first few
singular directions and decays rapidly outside the spike subspace.
We note that the strongest alignment need not occur at $k=1$, reflecting the
fact that the spike occupies a low-rank subspace rather than a single direction.

\begin{figure}[t]
  \centering
  \begin{subfigure}[t]{0.32\linewidth}
    \centering
    \includegraphics[width=\linewidth]{figures/col_batch.png}
    \caption{Batch mean.}
    \label{fig:col_batch}
  \end{subfigure}
  \hfill
  \begin{subfigure}[t]{0.32\linewidth}
    \centering
    \includegraphics[width=\linewidth]{figures/col_sampled.png}
    \caption{Sampled tokens across sequences.}
    \label{fig:col_sampled}
  \end{subfigure}
  \hfill
  \begin{subfigure}[t]{0.32\linewidth}
    \centering
    \includegraphics[width=\linewidth]{figures/col_single.png}
    \caption{Single-sequence mean.}
    \label{fig:col_single}
  \end{subfigure}

  \caption{Alignment between the broadcasted column-mean matrix and individual
  rank-$1$ singular components under different mean estimators.
  The horizontal axis is shown on a logarithmic scale.
  Alignment is sharply localized to the leading singular components, indicating
  that the spectral spike corresponds to a coherent, shared structure.}
  \label{fig:colmean_alignment}
\end{figure}

\paragraph{Concentration of alignment mass.}
While the alignment curves reveal the shape and decay of coherence along the
spectrum, they do not directly quantify how much of the shared structure is
captured by the spike subspace.
To make this explicit, we compute the cumulative distribution function (CDF) of
alignment mass, defined as the fraction of total alignment captured by the first
$k$ singular components.

Figure~\ref{fig:colmean_alignment_cdf_individual} shows the CDFs for each mean
estimator separately.
In all cases, alignment mass is overwhelmingly concentrated in the spike
subspace.
For the batch mean estimator, more than $99.6\%$ of the alignment mass is
contained within the top $1\%$ of singular components.
The sampled-token estimator exhibits a similarly sharp concentration
($98.7\%$ within the top $1\%$), while the single-sequence mean is more diffuse
but still places over $98\%$ of the alignment mass in the same spike subspace.

\begin{figure}[t]
  \centering
  \begin{subfigure}[t]{0.32\linewidth}
    \centering
    \includegraphics[width=\linewidth]{figures/col_batch_cdf.png}
    \caption{Batch mean.}
    \label{fig:col_batch_cdf}
  \end{subfigure}
  \hfill
  \begin{subfigure}[t]{0.32\linewidth}
    \centering
    \includegraphics[width=\linewidth]{figures/col_sampled_cdf.png}
    \caption{Sampled tokens across sequences.}
    \label{fig:col_sampled_cdf}
  \end{subfigure}
  \hfill
  \begin{subfigure}[t]{0.32\linewidth}
    \centering
    \includegraphics[width=\linewidth]{figures/col_single_cdf.png}
    \caption{Single-sequence mean.}
    \label{fig:col_single_cdf}
  \end{subfigure}

  \caption{CDF of alignment mass for individual mean estimators.
  Each plot shows the fraction of total alignment captured by the top $k$
  singular components.
  In all cases, alignment mass is overwhelmingly concentrated in the spike
  subspace, with the batch mean exhibiting the strongest concentration and the
  single-sequence mean the most diffuse distribution.}
  \label{fig:colmean_alignment_cdf_individual}
\end{figure}

To facilitate direct comparison across estimators, Figure~\ref{fig:colmean_alignment_cdf}
overlays the three CDFs in a single plot.
This view highlights the consistent ordering:
batch mean $>$ sampled tokens $>$ single sequence, reflecting the role of
semantic diversity in reinforcing shared structure under averaging.

\begin{figure}[t]
  \centering
  \includegraphics[width=0.7\linewidth]{figures/colmean_alignment_cdf_three_methods.png}
  \caption{Comparison of alignment mass concentration across mean estimators.
  The batch mean captures over $99.6\%$ of the total alignment mass within the
  top $1\%$ of singular components, followed by the sampled-token estimator
  ($98.7\%$) and the single-sequence mean ($98.2\%$).
  This ordering reflects the role of semantic diversity in reinforcing shared
  structure under averaging.}
  \label{fig:colmean_alignment_cdf}
\end{figure}

\paragraph{Effect of averaging across sequences.}
Taken together, Figures~\ref{fig:colmean_alignment}--\ref{fig:colmean_alignment_cdf}
establish that the spectral spike is not only low-rank, but also coherent and
globally shared across tokens.
Averaging over the full batch yields the strongest and most concentrated
alignment, followed by averaging over randomly sampled tokens across sequences,
while single-sequence averaging is least effective.
This behavior directly supports the view that the spike corresponds to shared
linguistic structure, rather than uncorrelated variance, and that its coherence
is precisely what enables simple averaging to isolate it.

This property distinguishes the spike from the high-dimensional, incoherent long
tail examined in the next subsection.

\subsection{The Long Tail: High-Dimensional and Incoherent}
\label{subsec:tail_incoherent}

We next analyze the long tail of the spectrum and show that it exhibits the
complementary properties needed for safe spike suppression:
(i) semantic information is distributed over a high-dimensional subspace, yet
each token occupies only a portion of this subspace (compressibility); and
(ii) semantic components are largely incoherent across tokens drawn from
different sequences or sufficiently distant positions, so averaging does not
erase semantic variation.

\paragraph{Per-token compressibility in the tail eigenbasis.}
We define the spike subspace as the span of the top $k_s$ right singular vectors
(e.g., $k_s=1\%$) and project each token onto the tail complement.
We then express each projected token in the \emph{tail eigenbasis} (the remaining
right singular vectors) and measure how many tail coordinates are required to
capture a fixed fraction of its energy.
Figure~\ref{fig:tail_effective_support} shows the resulting distributions of
effective support sizes.

In our running example, the tail subspace has dimension $H-k_s=4055$.
A typical token requires only a few hundred tail coordinates to capture
$90\%$ of its tail energy (median $\approx 542$), and under $1000$ coordinates to
capture $95\%$ (median $\approx 896$).
Thus, although the semantic subspace itself is high-dimensional, individual
tokens are \emph{compressible} in this basis: each token activates only a subset
of the available tail directions.

\paragraph{Coefficient-space incoherence across tokens.}
Compressibility alone does not guarantee safe averaging: if different tokens
activate the same directions with aligned coefficients, semantic information
could still be suppressed.
What matters is therefore not overlap of supports, but \emph{incoherence of
coefficients} in the tail basis.

To test this directly, we compute cosine similarity between normalized tail
coefficient vectors for pairs of tokens, distinguishing three cases:
(i) tokens from the same sequence and nearby positions;
(ii) tokens from the same sequence but far apart in position; and
(iii) tokens drawn from different sequences.
Figure~\ref{fig:tail_coeff_cosine} shows the resulting distributions.

Across sequences, tail coefficient cosine similarity is tightly concentrated
around zero (mean $\approx -1.1\times 10^{-3}$, standard deviation $\approx 0.05$),
indicating near-orthogonality in coefficient space.
Tokens from the same sequence but far apart in position exhibit a similar
distribution (mean $\approx -4.4\times 10^{-3}$), suggesting that semantic
correlation decays with distance in long contexts.
In contrast, nearby tokens within the same sequence show a modest positive shift
(mean $\approx 0.076$), reflecting local semantic and topical coherence.

\begin{figure}[t]
  \centering
  \begin{subfigure}[t]{0.48\linewidth}
    \centering
    \includegraphics[width=\linewidth]{figures/tail_effective_support.png}
    \caption{Per-token effective support size in the tail eigenbasis.}
    \label{fig:tail_effective_support}
  \end{subfigure}
  \hfill
  \begin{subfigure}[t]{0.48\linewidth}
    \centering
    \includegraphics[width=\linewidth]{figures/tail_coeff_cosine.png}
    \caption{Cosine similarity of tail coefficients:
    within-sequence (near), within-sequence (far), and across-sequence.}
    \label{fig:tail_coeff_cosine}
  \end{subfigure}
  \caption{Empirical properties of the long tail.
  (\textbf{Left}) Individual tokens are compressible in the semantic tail: most
  tail energy is concentrated on a subset of directions.
  (\textbf{Right}) Tail coefficients are nearly orthogonal across sequences and
  for far-apart positions within a sequence, while nearby tokens show modest
  correlation, indicating local semantic structure but global incoherence.}
\end{figure}

Taken together, these results complement
Section~\ref{subsec:spike_coherent}.
The spike is low-rank and coherent, corresponding to shared linguistic
structure, whereas the long tail is high-dimensional and semantically diverse:
individual tokens occupy only a portion of the tail subspace, and tail
coefficients are weakly correlated across sequences.
This separation provides the empirical basis for the analysis in the next
section, where we formalize why averaging can suppress the coherent spike while
preserving semantic information in the long tail.

\subsection{Elementwise Origin of Outliers: Mean vs Spike vs Tail}
\label{subsec:outlier_origin}

The previous subsections characterized anisotropy in spectral space.
However, model training operates in \emph{element space}.
To understand why averaging is effective, we must identify which component
(mean, spike, or tail) actually produces elementwise outliers.

We decompose the activation matrix $X$ into three orthogonal components:
\[
X = M + X_{\text{spike}} + X_{\text{tail}},
\]
where:
\begin{itemize}
    \item $M = \mathbf{1}\mu^\top$ is the batch-mean component,
    \item $X_{\text{spike}}$ is the projection onto the top-$k$ singular vectors
          (e.g., top $1\%$),
    \item $X_{\text{tail}}$ is the remaining orthogonal complement.
\end{itemize}

We then analyze the largest $0.1\%$ of element magnitudes (quantization-relevant outliers)
and compute, for each outlier entry, the squared-magnitude contribution fraction from
mean, spike, and tail components.

\paragraph{Outlier-level contribution.}

Figure~\ref{fig:outlier_contribution} shows the distribution of squared-magnitude
shares within the outlier set.

\begin{itemize}
    \item For \textbf{activations}, the overwhelming majority of outliers are dominated
          by the \emph{mean component}. Over $99\%$ of extreme entries have their largest
          squared contribution coming from $M$, while spike contributes modestly and tail
          contributes negligibly.
    \item For \textbf{gradient outputs}, the mean contribution is negligible.
          Outliers are instead dominated by variance components (primarily tail).
\end{itemize}

This demonstrates a fundamental asymmetry:
activation outliers are primarily a \emph{mean-driven phenomenon},
whereas gradient-output outliers are variance-driven.

\begin{figure}[t]
    \centering
    \includegraphics[width=0.95\linewidth]{figures/outlier_contribution_hist.png}
    \caption{
    Squared-magnitude contribution shares for the top $0.1\%$ elementwise outliers.
    (Left) Activation matrix.
    (Right) Gradient-output matrix.
    Activation outliers are dominated by the mean component,
    while gradient-output outliers arise from variance (tail) components.
    }
    \label{fig:outlier_contribution}
\end{figure}

\paragraph{Global energy decomposition.}

While outliers are dominated by the mean, the overall Frobenius energy tells a
different story.
Figure~\ref{fig:energy_decomposition} shows the fraction of total energy
attributed to each component.

\begin{itemize}
    \item For activations, the mean and spike each contribute roughly
          $20\%$ of total energy, while the long tail contributes the majority.
    \item For gradient outputs, the mean contributes almost no energy;
          the spectrum is dominated by spike and tail variance.
\end{itemize}

Thus:
\begin{itemize}
    \item The \emph{mean dominates extreme elements} (outliers).
    \item The \emph{tail dominates total variance}.
\end{itemize}

This distinction explains why mean removal dramatically reduces outliers
while preserving most variance.

\begin{figure}[t]
    \centering
    \includegraphics[width=0.75\linewidth]{figures/energy_decomposition.png}
    \caption{
    Global Frobenius energy decomposition into mean, spike, and tail components.
    Activation anisotropy is driven by a mean component that creates
    elementwise outliers, whereas gradient-output anisotropy is primarily variance-based.
    }
    \label{fig:energy_decomposition}
\end{figure}

\paragraph{Key implication.}

The results clarify the mechanism underlying averaging-based conditioning:

\begin{itemize}
    \item Mean removal eliminates the dominant source of activation outliers.
    \item It does not significantly disturb the long-tail variance.
    \item Therefore, it improves quantization stability without destroying semantic information.
\end{itemize}

This elementwise perspective resolves the earlier confusion:
anisotropy in spectral space does not imply that spike variance dominates
element magnitudes.
Instead, coherent mean structure is responsible for most extreme values,

\subsection{Activation Bias and Optimizer Second-Moment Inflation}
\label{subsec:optimizer_bias}

We now quantify how activation bias affects adaptive optimizers.
Adam and AdamW maintain an elementwise second-moment estimate
\[
v \approx \mathbb{E}[g_t \odot g_t],
\]
where for a linear layer
\[
g_t = \delta_t x_t^\top.
\]
Using the identity
\[
v = \frac{1}{N}(\Delta^{\odot 2})^\top (X^{\odot 2}),
\]
we can directly measure how conditioning the activations $X$
affects the distribution of $v$.

\paragraph{Protocol.}
For each layer and training step, we compute:
\begin{itemize}
\item $v_{\text{raw}}$ from original activations,
\item $v_{\text{mean}}$ after mean removal,
\item $v_{\text{mean+whiten}}$ after mean removal and top-$1\%$ spike whitening.
\end{itemize}
We focus on extreme-tail statistics ($p99.9$ and $\max$), which dominate
effective learning-rate suppression in Adam updates.

\paragraph{Quantitative results.}

Table~\ref{tab:adam_v_detailed} reports raw values and reduction factors.

\begin{table}[t]
\centering
\caption{Extreme-tail statistics of Adam second moment $v$.}
\label{tab:adam_v_detailed}
\begin{tabular}{lcccccc}
\toprule
Layer & Step & Method & $p99.9(v)$ & $\max(v)$ &
$p99.9$ reduction & $\max$ reduction \\
\midrule
\multirow{3}{*}{Layer 3}
& 10k
& Raw
& $1.60\times10^{-10}$ & $8.38\times10^{-10}$ & -- & -- \\
& 
& Mean-only
& $3.44\times10^{-11}$ & $1.91\times10^{-10}$ &
$\mathbf{4.6\times}$ & $\mathbf{4.4\times}$ \\
&
& Mean+whiten
& $1.72\times10^{-11}$ & $8.09\times10^{-11}$ &
$9.3\times$ & $10.4\times$ \\
\midrule
\multirow{3}{*}{Layer 3}
& 180k
& Raw
& $1.00\times10^{-10}$ & $7.37\times10^{-9}$ & -- & -- \\
&
& Mean-only
& $2.50\times10^{-11}$ & $1.53\times10^{-10}$ &
$\mathbf{4.0\times}$ & $\mathbf{48\times}$ \\
&
& Mean+whiten
& $2.44\times10^{-11}$ & $1.54\times10^{-10}$ &
$4.1\times$ & $47.8\times$ \\
\midrule
\multirow{3}{*}{Layer 21}
& 10k
& Raw
& $1.67\times10^{-11}$ & $3.43\times10^{-11}$ & -- & -- \\
&
& Mean-only
& $1.59\times10^{-11}$ & $2.58\times10^{-11}$ &
$1.05\times$ & $1.33\times$ \\
&
& Mean+whiten
& $1.48\times10^{-11}$ & $2.33\times10^{-11}$ &
$1.13\times$ & $1.47\times$ \\
\midrule
\multirow{3}{*}{Layer 21}
& 180k
& Raw
& $9.23\times10^{-11}$ & $1.10\times10^{-9}$ & -- & -- \\
&
& Mean-only
& $2.21\times10^{-11}$ & $2.69\times10^{-10}$ &
$\mathbf{4.2\times}$ & $\mathbf{4.1\times}$ \\
&
& Mean+whiten
& $2.09\times10^{-11}$ & $1.95\times10^{-10}$ &
$4.4\times$ & $5.6\times$ \\
\bottomrule
\end{tabular}
\end{table}

\paragraph{Observations.}

\begin{enumerate}
\item Mean removal consistently reduces the extreme tail of $v$.
In early layers or later training stages, reductions range from
$4\times$ to nearly $50\times$ in $\max(v)$.

\item Whitening provides additional improvement in early training
(Layer 3, 10k), but contributes little beyond mean removal
in later stages (Layer 3, 180k).

\item The median of $v$ changes minimally across methods,
indicating that mean removal primarily affects a small subset
of extreme parameters, which nonetheless dominate adaptive
learning-rate suppression.
\end{enumerate}

\paragraph{Implication.}
These results demonstrate that activation bias is a dominant driver of
extreme second-moment inflation in Adam.
Removing the mean significantly stabilizes the optimizer by shrinking
the pathological tail of $v$, thereby improving effective learning-rate
allocation without altering the semantic covariance structure.

% \section{A Structural Model of Linguistic and Semantic Components}
\label{sec:model}

Motivated by the empirical observations in Section~2, we introduce a structural
model that captures the essential geometric properties of language model
representations relevant to anisotropy and spike suppression.
The purpose of this model is not to reproduce all details of neural
representations, but to isolate the minimal structure needed to explain the
empirical behavior of averaging observed later in this section.
A formal explanation of these behaviors is provided in Section~4.

\subsection{Token Representation Model}
\label{subsec:token_model}

Let $x_i \in \mathbb{R}^H$ denote the representation of token $i$ at a fixed layer,
and let
\[
X = 
\begin{bmatrix}
x_1^\top \\
\vdots \\
x_N^\top
\end{bmatrix}
\in \mathbb{R}^{N \times H}
\]
collect the representations of all tokens in a batch.
We model each token representation as the sum of two components,
\begin{equation}
\label{eq:model_decomposition}
x_i = \ell + s_i,
\end{equation}
where:
\begin{itemize}
  \item $\ell \in \mathbb{R}^H$ is a \emph{shared linguistic component}, common to
  all tokens in the batch; and
  \item $s_i \in \mathbb{R}^H$ is a \emph{semantic component} that varies across
  tokens.
\end{itemize}

This decomposition reflects the empirical findings in
Sections~\ref{subsec:spike_long_tail}--\ref{subsec:tail_incoherent}:
linguistic structure (e.g., grammar, syntax, and shared functional patterns) is
globally coherent across tokens, while semantic content is diverse and
context-dependent.

\subsection{Structural Assumptions}
\label{subsec:structural_assumptions}

The empirical analyses in Sections~2.2--2.4 motivate the following structural
assumptions.

\paragraph{Assumption 1 (Low-rank linguistic structure).}
The shared component $\ell$ lies in a low-dimensional subspace
$\mathcal{L} \subset \mathbb{R}^H$ of dimension $d_\ell \ll H$.
Equivalently, $\ell$ is well approximated by the span of a small number of leading
right singular vectors of $X$, consistent with the spike structure observed in
Figure~\ref{fig:spectrum_spike_tail_combined} and the coherence analysis in
Figure~\ref{fig:colmean_alignment}.

\paragraph{Assumption 2 (Semantic compressibility).}
Each semantic component $s_i$ lies predominantly in the orthogonal complement
$\mathcal{L}^\perp$ and is \emph{compressible} in a fixed semantic eigenbasis:
most of its energy is concentrated on a subset of coordinates whose size is small
relative to $\dim(\mathcal{L}^\perp)$.
This assumption is empirically supported by the effective support statistics
shown in Figure~\ref{fig:tail_effective_support}.

\paragraph{Assumption 3 (Semantic incoherence across sequences).}
For tokens drawn from different sequences, or from sufficiently distant
positions within a long sequence, semantic components are weakly correlated in a
normalized sense.
Let $s_i^{\perp} = P_{\mathcal{L}^\perp}s_i$ denote the projection of $s_i$ onto the
semantic subspace.
Then for $i \neq j$ from different sequences,
\begin{equation}
\label{eq:assump_cos_incoh}
\mathbb{E}\!\left[
\frac{\langle s_i^{\perp},\, s_j^{\perp}\rangle}
{\|s_i^{\perp}\|_2\,\|s_j^{\perp}\|_2}
\right]
\approx 0,
\end{equation}
with bounded second moments $\mathbb{E}\|s_i\|_2^2 \le \sigma^2$.
This assumption is directly validated by the coefficient-space cosine similarity
distributions in Figure~\ref{fig:tail_coeff_cosine}.

\subsection{Empirical Effect of Averaging on Spike and Tail}
\label{subsec:avg_empirical}

We now empirically examine how mean removal (averaging) affects the spectral
structure characterized in Section~2.
The goal of this subsection is purely observational: to quantify how much of the
spectral spike is suppressed by averaging, and how much of the long tail is
preserved, without appealing to any theoretical explanation.
A formal account of these phenomena is deferred to Section~4.

\paragraph{Experimental setup.}
Given the token matrix $X \in \mathbb{R}^{N\times H}$, we consider two operations:
(i) \emph{mean removal},
\[
X_{\mathrm{avg}} = X - \mathbf{1}\mu^\top,
\quad
\mu = \tfrac{1}{N}\sum_{i=1}^N x_i,
\]
and (ii) \emph{top-$k$ SVD removal},
\[
X_k = X - X P_k,
\]
where $P_k$ projects onto the top-$k$ right singular vectors of $X$.
We compare these operations using Frobenius-norm energy ratios, which measure how
much variance remains after each transformation.

\paragraph{Mean alignment with the spike subspace.}
Figure~\ref{fig:avg_mean_in_spike} reports the fraction of mean energy captured by
the spike subspace, $\|P_k\mu\|_2^2 / \|\mu\|_2^2$, as a function of $k$.
The batch mean lies almost entirely within the spike subspace for modest $k$; for
example, projecting onto the top $1\%$ of singular directions captures over
$99.9\%$ of the mean energy.
This indicates that mean removal primarily targets directions aligned with the
dominant spectral spike.

\paragraph{Spike energy reduction and tail preservation.}
Figures~\ref{fig:avg_spike_reduction} and~\ref{fig:avg_tail_preservation} quantify
the effect of averaging on the spike and tail components, respectively.
Mean removal suppresses a nontrivial portion of spike energy, while preserving
nearly all tail energy across a wide range of spike-subspace definitions.
In contrast, top-$k$ SVD removal progressively eliminates both spike and tail
energy as $k$ increases.

\paragraph{Conservative nature of averaging.}
Figure~\ref{fig:avg_residual_energy} compares the total residual energy after mean
removal and after top-$k$ SVD removal.
Mean removal eliminates a fixed shared component, resulting in a bounded energy
reduction, whereas SVD removal continues to remove energy as the projection rank
grows.
This highlights the conservative nature of averaging: it suppresses the dominant
shared structure without aggressively removing variance from the long tail.

\begin{figure}[t]
  \centering
  \begin{subfigure}[t]{0.48\linewidth}
    \centering
    \includegraphics[width=\linewidth]{figures/avg_mean_in_spike_vs_k.png}
    \caption{Mean energy captured by the spike subspace.}
    \label{fig:avg_mean_in_spike}
  \end{subfigure}
  \hfill
  \begin{subfigure}[t]{0.48\linewidth}
    \centering
    \includegraphics[width=\linewidth]{figures/avg_spike_reduction_vs_k.png}
    \caption{Spike energy reduction under mean removal.}
    \label{fig:avg_spike_reduction}
  \end{subfigure}

  \vspace{0.5em}

  \begin{subfigure}[t]{0.48\linewidth}
    \centering
    \includegraphics[width=\linewidth]{figures/avg_tail_preservation_vs_k.png}
    \caption{Tail energy preservation under mean removal.}
    \label{fig:avg_tail_preservation}
  \end{subfigure}
  \hfill
  \begin{subfigure}[t]{0.48\linewidth}
    \centering
    \includegraphics[width=\linewidth]{figures/avg_residual_energy_comparison.png}
    \caption{Residual energy: mean removal vs.\ top-$k$ SVD removal.}
    \label{fig:avg_residual_energy}
  \end{subfigure}

  \caption{Empirical validation of averaging as conservative spike suppression.
  Mean removal targets directions aligned with the spike subspace
  (top-left), suppresses spike energy (top-right), preserves long-tail energy
  (bottom-left), and removes less total energy than top-$k$ SVD projection
  (bottom-right).}
  \label{fig:avg_vs_svd_validation}
\end{figure}

\paragraph{Summary of empirical findings.}
Across layers and batches, we consistently observe that mean removal:
(i) targets directions aligned with the spectral spike;
(ii) suppresses a substantial fraction of spike energy; and
(iii) preserves the vast majority of long-tail energy.
These empirical regularities set the stage for the next section, where we
provide a formal explanation for why averaging exhibits these properties under
the structural assumptions validated above.

% \section{Averaging as a Spectral Approximation}
\label{sec:averaging_spectral}

We now formalize the central claim of this work: under the anisotropic structure
empirically validated in Section~2 and modeled in Section~3, simple averaging
(mean subtraction) provides an effective and conservative approximation to
spectral spike suppression.
In particular, averaging removes a coherent component contained in the spike
subspace while preserving the semantic long tail up to vanishing error.

\subsection{Setup and Notation}

Let $X \in \mathbb{R}^{N\times H}$ denote the matrix of token representations in a
batch, with rows $x_i^\top$.
Recall the decomposition introduced in Section~\ref{sec:model}:
\[
x_i = \ell + s_i,
\qquad
X = \mathbf{1}\ell^\top + S,
\]
where $\ell$ is the shared linguistic component and $S$ collects the semantic
components.

Let $X = U\Sigma V^\top$ be the singular value decomposition of $X$, and let
$P_k = V_{1:k}^\top V_{1:k}$ denote the orthogonal projector onto the top-$k$
right singular vectors (the empirical \emph{spike subspace}).
Define the complementary projector $I-P_k$ as the \emph{tail subspace}.

Let
\[
\mu = \frac{1}{N}\sum_{i=1}^N x_i,
\qquad
X_{\mathrm{avg}} = X - \mathbf{1}\mu^\top
\]
denote the batch mean and the mean-removed matrix.

\subsection{Main Result}

\begin{theorem}[Averaging as conservative spike suppression]
\label{thm:avg_spectral}
Suppose the token representations satisfy Assumptions~1–3 of
Section~\ref{sec:model}, and that the empirical spike subspace $P_k$ contains the
linguistic subspace $\mathcal L$ up to small error (as observed empirically in
Section~2.3).
Then the averaging operator $X_{\mathrm{avg}}$ satisfies the following properties.

\begin{enumerate}

\item \textbf{Mean alignment with the spike subspace.}
The batch mean lies almost entirely in the empirical spike subspace:
\begin{equation}
\label{eq:mean_spike_align}
\frac{\|P_k \mu\|_2^2}{\|\mu\|_2^2} \ge 1 - \varepsilon,
\end{equation}
for modest $k$ and small $\varepsilon$.
Consequently, mean removal suppresses a coherent direction contained in the
spectral spike.
Empirically, this condition holds with $\varepsilon \ll 1$ for $k$ on the order
of the observed spike rank (Section~2.3).

\item \textbf{Spike energy reduction (conservative).}
Averaging removes a nonzero fraction of the spike energy:
\begin{equation}
\label{eq:spike_reduction}
1 - \frac{\|X_{\mathrm{avg}} P_k\|_F^2}{\|X P_k\|_F^2} > 0.
\end{equation}
Because averaging removes a single rank-one component aligned with the batch
mean, the relative spike reduction is largest when the spike subspace is small
and coherent.
As $k$ increases, the same removed direction constitutes a smaller fraction of
the larger spike subspace, and the relative reduction decreases.
This behavior reflects the conservative nature of averaging and is consistent
with the empirical curves in Section~3.

\item \textbf{Tail preservation (additive form).}
Mean removal perturbs the tail component only through the portion of the mean
lying outside the spike subspace:
\begin{equation}
\label{eq:tail_additive_main}
\|X_{\mathrm{avg}}(I-P_k) - X(I-P_k)\|_F
=
\|\mathbf{1}\mu^\top(I-P_k)\|_F
=
\sqrt{N}\,\|(I-P_k)\mu\|_2.
\end{equation}
In particular, when $\mu$ is well captured by the spike subspace (as in
Section~2.3), the additive perturbation to the semantic tail is small.

\smallskip
\noindent\emph{(Relative form under a tail-energy lower bound.)}
If, in addition, the tail energy satisfies $\|X(I-P_k)\|_F^2 \ge cN$ for some
constant $c>0$ (i.e., per-token tail variance is bounded below), then
\begin{equation}
\label{eq:tail_ratio_cor}
\frac{\|X_{\mathrm{avg}}(I-P_k)\|_F^2}{\|X(I-P_k)\|_F^2}
=
1 + O\!\left(\frac{\|(I-P_k)\mu\|_2}{\sqrt{N}}\right),
\end{equation}
so tail energy is preserved up to vanishing relative error as the number of
semantically independent tokens increases.

\end{enumerate}
\end{theorem}

\paragraph{Connection to gradient updates.}
Although the analysis above is stated in terms of token representations, the same
mean–residual separation applies directly to gradient computation in transformer
layers.
In particular, for linear projections, the batch gradient decomposes into a
shared rank-one mean–mean term and a cross-covariance term between error signals
and token-wise variations.
Mean removal eliminates only the shared rank-one component aligned with the
linguistic spike, while preserving the cross-covariance term that carries
token-specific learning signal.
As a result, averaging suppresses harmful coherent structure in gradients without
erasing semantic learning.
A detailed derivation of this gradient decomposition is provided in
Appendix~\ref{app:grad_decomposition}.

\subsection{Interpretation and Comparison with SVD}

Theorem~\ref{thm:avg_spectral} highlights a key distinction between averaging and
SVD-based spike suppression.
Top-$k$ SVD removal eliminates the entire spike subspace and removes increasing
energy as $k$ grows.
In contrast, averaging removes a single coherent direction aligned with the
batch mean.
It therefore provides a conservative approximation: it suppresses the most
harmful shared structure while minimizing collateral removal of semantic
variation.

This distinction explains the empirical curves in Section~3:
averaging removes a substantial fraction of spike energy while preserving almost
all tail energy, whereas SVD-based removal trades off spike suppression against
semantic loss as $k$ increases.

In the next section, we leverage this result to analyze the impact of averaging
on low-bit quantization and adaptive optimization.

% \section{Applications: Quantization and Optimization}
\label{sec:applications}

We now demonstrate how averaging-based spike suppression can be used in practice
to address two major challenges induced by anisotropy in large language model
training: low-bit quantization and adaptive optimization.
In both settings, the dominant spectral spike is known to be harmful, and prior
solutions rely on explicit spectral operations.
Our analysis shows that simple averaging can fill the role of SVD-based spike
suppression, providing a lightweight and hardware-friendly alternative.

\subsection{Low-Bit Quantization}
\label{subsec:quantization}

\paragraph{Problem setting.}
In low-bit and mixed-precision training, quantization error is dominated by
outliers and large-magnitude directions.
When representations are anisotropic, the dominant spectral spike sets the
numerical scale for quantization, forcing the remaining long-tail directions
into a narrow dynamic range.
This results in severe precision loss for semantically meaningful components and
has been identified as a key bottleneck in FP4 and FP8 training.

Recent approaches such as Metis mitigate this issue by explicitly suppressing
spectral spikes using SVD-based projections.
While effective, these methods require expensive matrix decompositions and are
difficult to integrate efficiently into large-scale training pipelines.

\paragraph{Averaging as a quantization-friendly approximation.}
The analysis in Sections~3 and~4 shows that the batch mean $\mu$ lies almost
entirely in the spike subspace and that mean removal suppresses a coherent
direction responsible for a large fraction of spike energy.
At the same time, tail energy is preserved up to vanishing error.
This makes averaging a natural preconditioning step for quantization:
\begin{enumerate}
  \item subtract the batch mean to suppress the dominant spike;
  \item quantize the mean-removed representations;
  \item (optionally) add the mean back in higher precision if needed.
\end{enumerate}

Unlike SVD-based methods, averaging is rank-one, requires only reductions and
broadcasts, and maps efficiently to GPUs and accelerators.
Empirically, we observe that mean removal significantly reduces spike-dominated
scale effects while preserving long-tail resolution, closely mirroring the
behavior of spectral spike suppression.
This explains why averaging can serve as a practical surrogate for SVD-based
methods in low-bit training regimes.

\paragraph{Interpretation.}
From the perspective of Section~4, quantization is sensitive to the \emph{range}
of representations.
Averaging removes the shared mean direction that dominates this range, thereby
rebalancing dynamic range without collapsing semantic variation.
In this sense, averaging provides a minimal intervention that targets the root
cause of quantization instability: coherent anisotropy.

\subsection{Adaptive Optimization}
\label{subsec:optimization}

\paragraph{Problem setting.}
Adaptive optimizers such as Adam and AdamW normalize gradients by estimates of
second moments.
In the presence of a strong spectral spike, these estimates are dominated by
shared, low-rank components.
As a result, learning rates in long-tail directions are implicitly suppressed,
slowing semantic learning and reducing effective model capacity.

Recent spike-aware optimizers, including Spectra, explicitly identify and
attenuate dominant spectral components in gradient statistics.
While effective, these methods again rely on explicit spectral analysis or
structured approximations.

\paragraph{Averaging as spike-aware normalization.}
Our results show that the batch mean captures the dominant shared component in
both activations and gradients.
Subtracting this mean before forming second-moment statistics reduces the
influence of the spike on adaptive normalization:
\[
g_{\mathrm{avg}} = g - \mathbb{E}[g],
\]
where $g$ denotes per-token or per-batch gradients.
By Theorem~\ref{thm:avg_spectral}, this operation suppresses a coherent spike
direction while preserving the covariance of semantic residuals.

In effect, averaging performs a form of spike-aware normalization at negligible
cost.
It attenuates the shared component that would otherwise dominate second-moment
estimates, allowing adaptive optimizers to allocate learning rates more evenly
across semantic directions.

\paragraph{Relation to spike-aware optimizers.}
From a spectral perspective, averaging removes a rank-one approximation to the
spike subspace.
This is weaker than full spectral suppression but sufficient to address the most
harmful shared component.
Empirically, this conservative behavior aligns with the design philosophy of
Spectra: suppress dominant modes without amplifying noise or overcorrecting
tail directions.

Thus, averaging provides a simple mechanism that approximates the effect of
spike-aware optimizers, while retaining the stability and simplicity of standard
adaptive methods.

\subsection{Summary}

Across both quantization and optimization, anisotropy manifests as a coherent
spectral spike that degrades performance.
Our analysis shows that this spike can be isolated through simple averaging,
without explicit spectral decomposition.
Averaging therefore fills the role of SVD-based spike suppression in practice:
it is inexpensive, hardware-friendly, and preserves the semantic structure that
drives learning.

These applications illustrate the practical impact of the “blessing” aspect of
anisotropy: the same coherence that creates training pathologies also enables
efficient remedies when properly exploited.

% \section{Related Work}
\label{sec:related_work}

Our work is closely related to three lines of research:
(i) spectral anisotropy in language model representations,
(ii) low-bit quantization under anisotropic spectra,
and (iii) optimizer design under structured gradient geometry.
We review each in turn and clarify how our contribution differs.

\paragraph{Anisotropy and outlier dimensions in language models.}
A growing body of work has documented that contextual representations produced
by Transformer language models are highly anisotropic.
Early studies showed that word and sentence embeddings concentrate in narrow
cones under cosine similarity, reducing representational diversity
\citep{arora2017simple,mu2017allbutthetop,ethayarajh2019contextual}.
Subsequent analyses identified small sets of ``rogue'' or ``outlier'' dimensions
that dominate similarity computations and downstream behavior
\citep{timkey2021rogue,kovaleva2021bertbusters,puccetti2022frequency}.
More recent work argues that anisotropy is closely tied to frequency skew,
syntactic structure, and training dynamics rather than being an inherent flaw
of the Transformer architecture \citep{rudman2023stable,machina2024anisotropy}.

Our work builds on these observations but shifts the focus from representation
quality at inference time to \emph{training-time spectral geometry}.
Rather than treating anisotropy as a nuisance to be removed, we show that the
coherent structure underlying anisotropy can be exploited algorithmically.

\paragraph{Low-bit quantization and spectral decomposition.}
Quantization of large language models has progressed rapidly from FP32 to BF16,
FP8, and recently FP4, motivated by substantial gains in efficiency
\citep{micikevicius2022fp8,peng2023fp8,perez2023fp8}.
However, extreme low-bit regimes exacerbate sensitivity to wide numerical
distributions and outliers.
A range of methods address this problem via reparameterization or preprocessing,
including channel-wise scaling and shifting \citep{xiao2023smoothquant,wei2023outlier},
orthogonal or Hadamard rotations \citep{suresh2017distributed,ashkboos2024quarot,ashkboos2025halo},
and outlier separation into high-precision branches \citep{li2024svdquant}.

Metis~\citep{cao2025metis} identifies spectral anisotropy as the \emph{fundamental}
barrier to FP4 training and proposes explicit spectral decomposition of weights,
activations, and gradients to separate dominant singular components from the
long tail.
By quantizing these components independently, Metis enables robust W4A4G4
training with minimal accuracy loss.
However, Metis still relies on estimating a low-rank spectral subspace, using
random sampling and projection to amortize the cost of SVD.

Our work is complementary but conceptually distinct.
Rather than approximating the spike subspace via randomized spectral methods,
we show that, due to the statistical structure of language data, a shared
linguistic component can be isolated by a much simpler operation—batch
averaging—without explicit decomposition.
This provides a lightweight surrogate that can fill the role of SVD-based spike
suppression in quantization-sensitive settings.

\paragraph{Optimizers under anisotropic gradient geometry.}
Adaptive optimizers such as Adam and AdamW \citep{kingma2014adam,loshchilov2017adamw}
operate coordinate-wise and ignore correlations within matrix-valued parameters.
When gradient energy concentrates in a low-rank spike, their second-moment
statistics become dominated by these directions, suppressing updates in the
long-tail subspace.
Matrix-aware optimizers, including K-FAC \citep{martens2015kfac}, Shampoo
\citep{gupta2018shampoo}, and SOAP \citep{vyas2024soap}, capture structured
correlations but incur significant computational and memory overhead.

Spectra~\citep{spectra2026} directly targets this issue by identifying and
suppressing the dominant low-rank spike subspace in gradient updates, while
leaving the noise-sensitive tail unamplified.
Spectra tracks the spike using cached power iteration and applies localized
spectral shaping with low overhead.

Our contribution is orthogonal to optimizer design.
We provide a data-geometric explanation for why the dominant spike is
\emph{coherent and shared across tokens}, and show that averaging removes a
direction lying almost entirely within the spike subspace.
This insight complements spike-aware optimizers such as Spectra by providing a
simple, hardware-friendly primitive that attenuates spike-dominated statistics
without requiring iterative spectral routines.

\paragraph{Efficient low-rank approximation.}
Classical SVD and QR-based methods are expensive in large-scale training loops.
Randomized and streaming techniques offer scalable alternatives for estimating
dominant subspaces \citep{halko2011randomized,tropp2012user,oja1982pca}.
Both Metis and Spectra exploit these tools to reduce the cost of spectral
interventions.

Our work shows that, under the specific statistical structure of human language
corpora, even these approximations may be unnecessary for isolating the most
harmful shared component.
By leveraging coherence in linguistic structure and incoherence in semantic
residuals, averaging acts as a conservative but effective spectral
approximation.

% \section{Discussion, Limitations, and Outlook}
\label{sec:discussion}

This work reframes anisotropy in human language representations as a dual
phenomenon.
On one hand, the presence of a dominant spectral spike poses significant
challenges for optimization and quantization.
On the other hand, the coherence and shared nature of this spike make it
amenable to simple, efficient suppression through averaging.
By grounding this perspective in empirical analysis, a minimal structural
model, and formal guarantees, we show that averaging can fill the role of
spectral spike suppression in a principled and hardware-friendly manner.

\subsection{When Averaging Works}

The effectiveness of averaging relies on structural properties that are
ubiquitous in large language model training, but not universal.
In particular, averaging is most effective when:
\begin{itemize}
  \item the dominant spike corresponds to shared linguistic structure that is
  coherent across tokens;
  \item semantic components are diverse and weakly correlated across sequences
  or sufficiently distant contexts;
  \item batches contain a sufficient number of semantically independent tokens
  so that semantic residuals cancel under averaging.
\end{itemize}

These conditions are empirically validated for standard language model training
regimes, but may not hold in settings with highly correlated batches, extremely
small batch sizes, or non-linguistic data distributions.
In such cases, more aggressive spectral methods may still be required.

\subsection{Conservatism versus Aggressiveness}

A key theme throughout the paper is that averaging is a \emph{conservative}
operation.
Unlike SVD-based spike suppression, which removes an entire low-rank subspace,
averaging removes only a single coherent direction aligned with the batch mean.
This conservatism is both a strength and a limitation.

On the positive side, it minimizes collateral damage to semantic structure and
avoids amplifying noise in the tail.
On the negative side, it cannot eliminate all spike energy when the spike
subspace has dimension greater than one.
As a result, averaging should be viewed as a safe first-order approximation,
rather than a complete replacement for spectral methods in all settings.

\subsection{Relation to Existing Spectral Methods}

Our analysis clarifies why explicit spectral methods such as Metis and
spike-aware optimizers are effective: they target the same coherent structure
identified here.
The contribution of this work is not to dispute the value of such methods, but
to show that much of their benefit can be achieved by exploiting statistical
regularities of language data through simple operations.

From this perspective, averaging can be seen as a baseline spectral
approximation.
It may be combined with more refined methods, used as a warm-up or
preconditioning step, or applied selectively in regimes where full spectral
analysis is impractical.

\subsection{Broader Implications}

The curse–blessing duality highlighted in this work suggests a broader principle
for large-scale learning systems:
pathological geometric structures induced by data may also encode exploitable
regularities.
Understanding the geometry of data distributions—not just model architectures—
can open the door to algorithmic simplifications that are both theoretically
sound and practically impactful.

While this paper focuses on language models, similar anisotropic structures have
been observed in vision, multimodal models, and scientific data.
Exploring whether analogous averaging-based approximations apply in these
domains is an interesting direction for future work.

\subsection{Conclusion}

Anisotropy in human language corpora is unavoidable.
It degrades training efficiency, destabilizes quantization, and complicates
optimization.
At the same time, its coherence enables simple and efficient remedies.

This work shows that averaging, one of the simplest operations available in
large-scale systems, can be elevated from a heuristic to a principled spectral
approximation.
By aligning algorithmic design with the geometry of language data, we obtain a
method that is efficient, interpretable, and broadly applicable.
We hope this perspective encourages further exploration of data-aware
approaches to large-scale model training.

\newpage


\bibliographystyle{../doc/iclr-template/master/iclr2026/iclr2026_conference}
\bibliography{references}

\newpage

\appendix
%\appendix
\section{Formal Proofs for Averaging as Spectral Approximation}
\label{app:proofs}

This appendix provides rigorous proofs supporting
Section~\ref{sec:averaging_spectral}.
To avoid circularity, we first prove results with respect to a fixed linguistic
subspace $\mathcal L$, and then relate it to the empirical SVD spike subspace.

\subsection{Setup}

Assume token representations satisfy
\[
x_i = \ell + s_i,
\qquad
X = \mathbf{1}\ell^\top + S,
\]
where $\ell \in \mathcal L \subset \mathbb{R}^H$ and $s_i$ are semantic residuals.

Let $P_{\mathcal L}$ and $P_{\mathcal L^\perp}$ denote orthogonal projectors onto
$\mathcal L$ and its complement.

\subsection{Assumptions}

\begin{assumption}[Low-rank linguistic structure]
$\ell \in \mathcal L$, where $\dim(\mathcal L)=d_\ell \ll H$.
\end{assumption}

\begin{assumption}[Semantic second-moment bound]
There exists $\sigma^2$ such that
\[
\mathbb{E}\|s_i\|_2^2 \le \sigma^2.
\]
\end{assumption}

\begin{assumption}[Semantic incoherence across sequences]
For tokens drawn from different sequences (or sufficiently distant positions),
the normalized semantic inner product is small. Specifically, for $i \neq j$,
\begin{equation}
\label{eq:assump_cos_incoh}
\left|
\mathbb{E}\!\left[
\frac{s_i^\top s_j}{\|s_i\|_2\,\|s_j\|_2}
\right]
\right|
\le \varepsilon_{\mathrm{inc}},
\end{equation}
for a small constant $\varepsilon_{\mathrm{inc}} \ll 1$.
Moreover, in the tail eigenbasis used in Section~2.4, the normalized coefficient
vectors are also weakly aligned across sequences, i.e.,
$\left|\mathbb{E}[\cos(c_i, c_j)]\right| \le \varepsilon_{\mathrm{inc}}$, where
$c_i$ denotes the tail-basis coefficient vector of $s_i$.
\end{assumption}

This assumption directly matches the empirical cosine similarity measurements in
Section~2.4.

\subsection{Mean Removal Identity}

\begin{lemma}
Let $\bar{s}_N = \frac{1}{N}\sum_{i=1}^N s_i$.
Then
\[
\mu = \ell + \bar{s}_N,
\qquad
X_{\mathrm{avg}} = S - \mathbf{1}\bar{s}_N^\top.
\]
\end{lemma}

\subsection{Concentration of Semantic Mean}

\begin{lemma}
With probability at least $1-\delta$,
\[
\|\bar{s}_N\|_2 \le \frac{\sigma}{\sqrt{N\delta}}.
\]
\end{lemma}

\subsection{Tail Preservation under Mean Removal}
\label{app:tail_preservation}

\begin{theorem}[Tail perturbation is rank-one and controlled by mean leakage]
\label{thm:app_tail_additive}
For any projector $P$ on $\mathbb{R}^H$ and any batch mean $\mu$,
\begin{equation}
\label{eq:app_tail_exact}
X_{\mathrm{avg}}(I-P) - X(I-P) = -\mathbf{1}\mu^\top(I-P),
\end{equation}
and therefore
\begin{equation}
\label{eq:app_tail_bound}
\|X_{\mathrm{avg}}(I-P) - X(I-P)\|_F
=
\|\mathbf{1}\mu^\top(I-P)\|_F
=
\sqrt{N}\,\|(I-P)\mu\|_2.
\end{equation}
In particular, when $P=P_k$ is the empirical spike projector and $\mu$ is well
captured by $\mathrm{range}(P_k)$, tail perturbation is small.
\end{theorem}

\begin{proof}
Equation~\eqref{eq:app_tail_exact} follows directly from
$X_{\mathrm{avg}} = X - \mathbf{1}\mu^\top$.
For \eqref{eq:app_tail_bound}, use the Frobenius norm identity
$\|\mathbf{1}w^\top\|_F = \|\mathbf{1}\|_2\|w\|_2 = \sqrt{N}\|w\|_2$ with
$w=(I-P)\mu$.
\end{proof}

\begin{corollary}[Relative tail preservation under a tail-energy lower bound]
\label{cor:app_tail_ratio}
Assume $\|X(I-P)\|_F^2 \ge cN$ for some constant $c>0$.
Then
\[
\frac{\|X_{\mathrm{avg}}(I-P)\|_F^2}{\|X(I-P)\|_F^2}
=
1 + O\!\left(\frac{\|(I-P)\mu\|_2}{\sqrt{N}}\right).
\]
\end{corollary}

\begin{proof}
Let $A=X(I-P)$ and $E=\mathbf{1}\mu^\top(I-P)$.
Then $X_{\mathrm{avg}}(I-P)=A-E$ and
\[
\bigl|\|A-E\|_F^2 - \|A\|_F^2\bigr|
\le 2\|A\|_F\|E\|_F + \|E\|_F^2.
\]
Divide both sides by $\|A\|_F^2$ and use $\|E\|_F=\sqrt{N}\|(I-P)\mu\|_2$ and
$\|A\|_F^2\ge cN$.
\end{proof}

\paragraph{Definition (Semantic leakage into the linguistic subspace).}
We say semantic leakage into $\mathcal L$ is small if
\[
\mathbb{E}\|P_{\mathcal L} s_i\|_2^2 \le \rho^2,
\qquad \rho \ll \|\ell\|_2.
\]

\subsection{Spike Suppression}

\begin{theorem}
Mean removal eliminates the coherent linguistic component:
\[
XP_{\mathcal L} = \mathbf{1}\ell^\top + SP_{\mathcal L},
\qquad
X_{\mathrm{avg}}P_{\mathcal L} = SP_{\mathcal L} - \mathbf{1}\bar{s}_N^\top P_{\mathcal L}.
\]
If semantic leakage into $\mathcal L$ is small in the above sense, then
\[
\|X_{\mathrm{avg}}P_{\mathcal L}\|_F^2 \ll \|XP_{\mathcal L}\|_F^2.
\]
\end{theorem}

\subsection{Relation to the Empirical SVD Spike Subspace}

Let $P_k$ denote the projector onto the top-$k$ right singular vectors of $X$.
Empirically (Section~2.3),
\[
\frac{\|P_k \mu\|_2^2}{\|\mu\|_2^2} \approx 1
\]
for modest $k$, indicating that $\mathcal L \subseteq \mathrm{range}(P_k)$ up to
small error.
Under this condition, the fixed-subspace guarantees above apply to the empirical
spike subspace used in practice.

\paragraph{Remark.}
Because averaging removes exactly one rank-one component, its relative spike
reduction decreases as $k$ grows.
This behavior matches the empirical spike-reduction curves in Section~3 and is a
direct consequence of the conservative nature of averaging.

\section{Gradient Decomposition Under Mean-Centered Activations}
\label{app:grad_decomposition}

This appendix clarifies how mean-centering activations interacts with gradient
computation in transformer layers, and why removing a shared mean component does
not eliminate token-specific learning signal, even though per-token gradients
are outer products.
The discussion here is explanatory and complements the formal analysis in
Section~4.

\paragraph{Per-token gradient structure.}
Consider a linear projection in a transformer layer,
\[
y = Wx, \qquad W \in \mathbb{R}^{d \times d}, \; x \in \mathbb{R}^d,
\]
and let $\ell_t$ denote the loss associated with token $t$.
The gradient of $\ell_t$ with respect to $W$ is
\[
g_t \triangleq \nabla_W \ell_t = \delta_t x_t^\top,
\]
where $\delta_t = \nabla_{y_t} \ell_t$ is the backpropagated error signal.
The batch gradient used by SGD is the average of these outer products,
\[
G = \frac{1}{N}\sum_{t=1}^N g_t = \mathbb{E}[\delta x^\top].
\]

\paragraph{Decomposition into mean and residual components.}
Decompose the activation into its mean and residual components,
\[
x = \mu_x + \tilde{x},
\qquad
\mu_x = \mathbb{E}[x],
\qquad
\mathbb{E}[\tilde{x}] = 0.
\]
Substituting into the batch gradient yields the exact identity
\begin{equation}
\label{eq:grad_decomp}
\mathbb{E}[\delta x^\top]
=
\underbrace{\mathbb{E}[\delta]\mu_x^\top}_{\text{rank-1, shared}}
+
\underbrace{\mathbb{E}[\delta \tilde{x}^\top]}_{\text{cross-covariance}}.
\end{equation}

\noindent
Equation~\eqref{eq:grad_decomp} separates the batch gradient into two conceptually
distinct components:
\begin{itemize}
  \item the mean term $\mathbb{E}[\delta]\mu_x^\top$, which is a shared, low-rank
  update arising from components that are coherent across tokens (the
  grammar-driven spike); and
  \item the cross-covariance term $\mathbb{E}[\delta \tilde{x}^\top]$, which
  captures correlations between error signals and token-wise variations and
  constitutes the learning-relevant semantic signal.
\end{itemize}

\paragraph{Effect of mean-centering.}
If activations are mean-centered, $x' = x - \mu_x = \tilde{x}$, then the batch
gradient becomes
\[
\mathbb{E}[\delta (x')^\top] = \mathbb{E}[\delta \tilde{x}^\top].
\]
Thus, mean-centering removes only the rank-one mean--mean contribution
$\mathbb{E}[\delta]\mu_x^\top$ and leaves the cross-covariance term intact.
This explains why averaging suppresses a coherent, spike-like component of the
gradient while preserving token-specific learning signal.

\paragraph{Relation to second-order structure and adaptive optimizers.}
Although SGD updates use the first moment of per-token gradients, the learning
signal associated with semantic variation is expressed through correlations
between $\delta_t$ and $\tilde{x}_t$ across tokens.
Moreover, adaptive optimizers such as Adam and AdamW explicitly maintain
second-moment estimates of gradients,
\[
v \approx \mathbb{E}[g_t \odot g_t],
\]
which decompose into contributions from shared components and from token-wise
variability.
A large shared component inflates these second moments and suppresses effective
learning rates in semantically meaningful directions.
Mean-centering reduces this shared contribution while leaving the variability
term largely unchanged, which is consistent with the empirical improvements
observed in optimization and low-bit quantization.

\paragraph{Summary.}
Although per-token gradients are outer products, mean-centering activations
removes only a shared, low-rank component of the batch gradient.
The cross-covariance term that carries token-specific semantic information
remains, ensuring that semantic learning is preserved.
This decomposition provides an intuitive bridge between the empirical results in
Section~3 and the formal analysis in Section~4.
% Extended derivations


\end{document}